%!TEX TS-program = lualatex
%!TEX encoding = UTF-8

\documentclass[12pt,a4paper,ja=standard]{bxjsarticle}

% 数学パッケージ
\usepackage{amsmath,amssymb,amsthm}
\usepackage{mathtools}

% 図表パッケージ
\usepackage{tikz}
\usepackage{tikz-cd}
\usetikzlibrary{arrows,positioning,shapes,calc,matrix}

% レイアウト
\usepackage[margin=2.5cm]{geometry}
\usepackage{enumitem}
\usepackage{tcolorbox}
\tcbuselibrary{theorems,skins,breakable}

% ハイパーリンク
\usepackage{hyperref}
\hypersetup{
    colorlinks=true,
    linkcolor=blue,
    filecolor=magenta,      
    urlcolor=cyan,
}

% 定理環境
\theoremstyle{definition}
\newtheorem{definition}{定義}[section]
\newtheorem{theorem}[definition]{定理}
\newtheorem{lemma}[definition]{補題}
\newtheorem{proposition}[definition]{命題}
\newtheorem{example}[definition]{例}
\newtheorem{remark}[definition]{注意}

% カスタムボックス
\newtcolorbox{keypoint}[1][]{
    colback=blue!5!white,
    colframe=blue!75!black,
    fonttitle=\bfseries,
    title=重要ポイント,
    #1
}

\newtcolorbox{insight}[1][]{
    colback=green!5!white,
    colframe=green!75!black,
    fonttitle=\bfseries,
    title=数学的洞察,
    #1
}

\newtcolorbox{connection}[1][]{
    colback=orange!5!white,
    colframe=orange!75!black,
    fonttitle=\bfseries,
    title=構造塔との接続,
    #1
}

% タイトル情報
\title{\Huge\bfseries 離散確率論の代数的構造\\
\Large DecidableAlgebra.lean: 有限代数と可測性}

\author{
    \large Lean 4 コード: Codex (OpenAI) \\
    \large \TeX 解説: Claude Code \\
    \vspace{0.3cm}
    \normalsize プロジェクト: 構造塔理論と確率論
}

\date{\today}

\begin{document}

\maketitle

\begin{abstract}
本稿では、\texttt{DecidableEvents.lean} の自然な拡張として、有限標本空間における\textbf{有限代数}（finite algebra）の形式化について解説する。有限代数は Boolean 演算で閉じた事象族であり、$\sigma$-代数の有限版に相当する。これにより、「観測可能な情報の構造」を数学的に記述でき、将来の DecidableFiltration（構造塔のインスタンス化）への準備が整う。本解説は学部 3--4 年生を対象とし、測度論の基礎知識を前提とする。
\end{abstract}

\tableofcontents
\newpage

\section{序論：事象から代数へ}

\subsection{前回のまとめ}

\texttt{DecidableEvents.lean} では、以下を学んだ：

\begin{itemize}
    \item 有限標本空間 $\Omega$ の定義
    \item 事象 $A \subseteq \Omega$ と decidability
    \item 基本演算（補集合、和、積）
    \item 具体例（サイコロ、コイン投げ）
\end{itemize}

しかし、個々の事象だけでは不十分である。確率論では、\textbf{事象の族}（family of events）を扱う必要がある。

\subsection{なぜ事象の族が必要か}

\begin{example}[情報の増加]
3回コイン投げの実験を考える。標本空間は $\Omega = \{H, T\}^3$（8通り）。

\begin{itemize}
    \item \textbf{時刻 0}：何も知らない $\Rightarrow$ 識別可能な事象：$\{\emptyset, \Omega\}$（2個）
    \item \textbf{時刻 1}：1回目の結果を知る $\Rightarrow$ 識別可能な事象：4個
    \item \textbf{時刻 2}：2回目まで知る $\Rightarrow$ 識別可能な事象：8個
    \item \textbf{時刻 3}：すべて知る $\Rightarrow$ 識別可能な事象：16個
\end{itemize}

各時刻で「観測可能な事象の族」が必要になる。
\end{example}

\subsection{本稿の位置づけ}

\begin{center}
\begin{tikzpicture}[
    box/.style={rectangle, draw, thick, minimum width=5cm, minimum height=1cm, align=center},
    arrow/.style={->, thick}
]
    \node[box] (level1) at (0,0) {1. DecidableEvents.lean\\{\small 事象の基礎}};
    \node[box, fill=blue!10, below=0.5cm of level1] (level2) {2. DecidableAlgebra.lean\\{\small 事象の構造化された族}};
    \node[box, below=0.5cm of level2] (level3) {3. DecidableFiltration.lean\\{\small 構造塔のインスタンス}};
    \node[box, below=0.5cm of level3] (level4) {4. ComputableStoppingTime.lean\\{\small 停止時間}};
    \node[box, below=0.5cm of level4] (level5) {5. AlgorithmicMartingale.lean\\{\small マルチンゲール理論}};
    
    \draw[arrow] (level1) -- (level2);
    \draw[arrow] (level2) -- (level3);
    \draw[arrow] (level3) -- (level4);
    \draw[arrow] (level4) -- (level5);
    
    \node[right=0.3cm of level2, text width=3cm, font=\small] {{\color{blue}$\leftarrow$ 今回}};
    \node[right=0.3cm of level3, text width=3cm, font=\small] {{\color{red}構造塔理論}};
\end{tikzpicture}
\end{center}

\subsection{本稿の構成}

\begin{enumerate}
    \item 第 \ref{sec:finite-algebra} 節：有限代数の定義と基本性質
    \item 第 \ref{sec:examples} 節：具体例（自明な代数、偶奇代数）
    \item 第 \ref{sec:measurable} 節：可測関数
    \item 第 \ref{sec:connection} 節：構造塔との接続
    \item 第 \ref{sec:conclusion} 節：まとめと展望
\end{enumerate}

\section{有限代数の定義}\label{sec:finite-algebra}

\subsection{$\sigma$-代数の復習}

\begin{definition}[$\sigma$-代数]
標本空間 $\Omega$ 上の事象族 $\mathcal{F} \subseteq \mathcal{P}(\Omega)$ が \textbf{$\sigma$-代数}（sigma-algebra）であるとは：

\begin{enumerate}
    \item \textbf{全事象}：$\Omega \in \mathcal{F}$
    \item \textbf{補集合で閉じている}：$A \in \mathcal{F} \Rightarrow A^c \in \mathcal{F}$
    \item \textbf{可算和で閉じている}：$(A_n)_{n \in \mathbb{N}} \subseteq \mathcal{F} \Rightarrow \bigcup_{n=1}^{\infty} A_n \in \mathcal{F}$
\end{enumerate}
\end{definition}

\begin{remark}
組 $(\Omega, \mathcal{F})$ を\textbf{可測空間}（measurable space）と呼ぶ。
\end{remark}

\subsection{有限代数の定義}

有限標本空間では、可算和の条件を\textbf{有限和}に弱めても十分である。

\begin{definition}[有限代数]
標本空間 $\Omega$ 上の事象族 $\mathcal{F} \subseteq \mathcal{P}(\Omega)$ が \textbf{有限代数}（finite algebra）であるとは：

\begin{enumerate}
    \item \textbf{空事象}：$\emptyset \in \mathcal{F}$（または同値に $\Omega \in \mathcal{F}$）
    \item \textbf{補集合で閉じている}：$A \in \mathcal{F} \Rightarrow A^c \in \mathcal{F}$
    \item \textbf{和で閉じている}：$A, B \in \mathcal{F} \Rightarrow A \cup B \in \mathcal{F}$
\end{enumerate}
\end{definition}

\begin{tcolorbox}[title=Lean コード（170--179 行目）]
\begin{verbatim}
structure FiniteAlgebra (Ω : Type*) where
  /-- 事象の族 -/
  events : Set (Event Ω)
  /-- 空事象を含む -/
  has_empty : Event.empty ∈ events
  /-- 補集合で閉じている -/
  closed_complement : ∀ {A : Event Ω}, 
    A ∈ events → Event.complement A ∈ events
  /-- 和で閉じている -/
  closed_union : ∀ {A B : Event Ω}, 
    A ∈ events → B ∈ events → Event.union A B ∈ events
\end{verbatim}
\end{tcolorbox}

\subsection{なぜ「代数」と呼ぶか}

\begin{insight}
有限代数は\textbf{代数的構造}を持つ：

\begin{center}
\begin{tabular}{|l|l|l|}
\hline
\textbf{演算} & \textbf{代数的役割} & \textbf{記号} \\
\hline
補集合 & 単項演算 & $(\cdot)^c : \mathcal{F} \to \mathcal{F}$ \\
和集合 & 二項演算 & $\cup : \mathcal{F} \times \mathcal{F} \to \mathcal{F}$ \\
積集合 & 二項演算 & $\cap : \mathcal{F} \times \mathcal{F} \to \mathcal{F}$ \\
\hline
\end{tabular}
\end{center}

これらの演算で\textbf{閉じている} = 代数的構造
\end{insight}

\subsection{基本的な性質}

\begin{theorem}[全事象の包含]
$\mathcal{F}$ が有限代数ならば、$\Omega \in \mathcal{F}$。
\end{theorem}

\begin{proof}
$\emptyset \in \mathcal{F}$ かつ補集合で閉じているので、
\[
\emptyset^c = \Omega \in \mathcal{F}
\]
\end{proof}

\begin{theorem}[積で閉じている]
$\mathcal{F}$ が有限代数で $A, B \in \mathcal{F}$ ならば、$A \cap B \in \mathcal{F}$。
\end{theorem}

\begin{proof}
De Morgan の法則を使用：
\[
A \cap B = (A^c \cup B^c)^c
\]

\begin{enumerate}
    \item $A, B \in \mathcal{F}$ より、$A^c, B^c \in \mathcal{F}$（補集合で閉じている）
    \item $A^c \cup B^c \in \mathcal{F}$（和で閉じている）
    \item $(A^c \cup B^c)^c \in \mathcal{F}$（補集合で閉じている）
\end{enumerate}
\end{proof}

\begin{center}
\begin{tikzpicture}
    % De Morgan の視覚化
    \draw[thick] (0,0) rectangle (6,3);
    \node at (3,3.3) {\textbf{De Morgan の法則}};
    
    % 左側：A ∩ B
    \begin{scope}
        \clip (1,1.5) circle (0.8);
        \fill[blue!30] (2,1.5) circle (0.8);
    \end{scope}
    \draw[thick] (1,1.5) circle (0.8);
    \draw[thick] (2,1.5) circle (0.8);
    \node at (0.5,2.5) {$A$};
    \node at (2.5,2.5) {$B$};
    \node at (1.5,0.3) {$A \cap B$};
    
    % 矢印
    \draw[->, very thick] (3,1.5) -- (3.5,1.5);
    \node at (3.25,1.8) {$=$};
    
    % 右側：(A^c ∪ B^c)^c
    \begin{scope}[xshift=3.5cm]
        \fill[blue!30] (0,0) rectangle (2.5,3);
        \fill[white] (0.5,1.5) circle (0.8);
        \fill[white] (1.5,1.5) circle (0.8);
        \draw[thick] (0.5,1.5) circle (0.8);
        \draw[thick] (1.5,1.5) circle (0.8);
        \node at (0,2.5) {$A^c$};
        \node at (2,2.5) {$B^c$};
        \node at (1,0.3) {$(A^c \cup B^c)^c$};
    \end{scope}
\end{tikzpicture}
\end{center}

\begin{theorem}[有限和・積で閉じている]
$\mathcal{F}$ が有限代数ならば：
\begin{enumerate}
    \item 任意の有限個の事象 $A_1, \ldots, A_n \in \mathcal{F}$ に対して、$\bigcup_{i=1}^n A_i \in \mathcal{F}$
    \item 任意の有限個の事象 $A_1, \ldots, A_n \in \mathcal{F}$ に対して、$\bigcap_{i=1}^n A_i \in \mathcal{F}$
\end{enumerate}
\end{theorem}

\begin{proof}
帰納法による。
\end{proof}

\section{具体例}\label{sec:examples}

\subsection{自明な代数}

\begin{definition}[自明な代数]
\textbf{自明な代数}（trivial algebra）は：
\[
\mathcal{F}_{\text{trivial}} = \{\emptyset, \Omega\}
\]

これは最小の代数であり、「何も観測できない」状態を表す。
\end{definition}

\begin{example}[時刻 0 の情報]
確率実験の開始前（時刻 0）では、まだ何も観測していない。この状態では：
\begin{itemize}
    \item $\Omega$ は「何が起こるか全く分からない」
    \item $\emptyset$ は「絶対に起こらない」
\end{itemize}
のみが識別可能。
\end{example}

\subsection{全体の代数}

\begin{definition}[全体の代数]
\textbf{全体の代数}（power set algebra）は：
\[
\mathcal{F}_{\text{power}} = \mathcal{P}(\Omega)
\]

これは最大の代数であり、「すべてを観測できる」状態を表す。
\end{definition}

\begin{example}[完全な情報]
すべての実験が終了した後では、すべての事象について「起こったか・起こらなかったか」が判定できる。
\end{example}

\subsection{偶奇代数（サイコロの例）}

\begin{example}[偶奇で生成される代数]
6面サイコロ $\Omega = \{0, 1, 2, 3, 4, 5\}$ に対して、偶数の目の事象：
\[
E = \{0, 2, 4\}
\]

から生成される代数は：
\[
\mathcal{F}_{\text{even-odd}} = \{\emptyset, E, E^c, \Omega\}
\]

ここで：
\begin{itemize}
    \item $\emptyset$ = 空事象
    \item $E = \{0, 2, 4\}$ = 偶数の目
    \item $E^c = \{1, 3, 5\}$ = 奇数の目
    \item $\Omega = \{0, 1, 2, 3, 4, 5\}$ = 全体
\end{itemize}
\end{example}

\begin{tcolorbox}[title=Lean コード（430--436 行目）]
\begin{verbatim}
def evenOddAlgebra : FiniteAlgebra diceSample.carrier where
  events := {
    Event.empty,
    evenDice,
    Event.complement evenDice,
    Event.univ
  }
  -- 代数の公理を証明
\end{verbatim}
\end{tcolorbox}

\subsubsection{偶奇代数の視覚化}

\begin{center}
\begin{tikzpicture}[scale=1.2]
    % サイコロの目
    \foreach \i in {0,...,5} {
        \pgfmathsetmacro{\x}{mod(\i,3)*2}
        \pgfmathsetmacro{\y}{-int(\i/3)*2}
        
        \ifnum\i=0
            \node[circle, draw, thick, fill=blue!20, minimum size=1cm] at (\x,\y) {\Large \i};
        \fi
        \ifnum\i=1
            \node[circle, draw, thick, fill=red!20, minimum size=1cm] at (\x,\y) {\Large \i};
        \fi
        \ifnum\i=2
            \node[circle, draw, thick, fill=blue!20, minimum size=1cm] at (\x,\y) {\Large \i};
        \fi
        \ifnum\i=3
            \node[circle, draw, thick, fill=red!20, minimum size=1cm] at (\x,\y) {\Large \i};
        \fi
        \ifnum\i=4
            \node[circle, draw, thick, fill=blue!20, minimum size=1cm] at (\x,\y) {\Large \i};
        \fi
        \ifnum\i=5
            \node[circle, draw, thick, fill=red!20, minimum size=1cm] at (\x,\y) {\Large \i};
        \fi
    }
    
    \node[blue] at (-1.5,0) {偶数 $E$};
    \node[red] at (-1.5,-2) {奇数 $E^c$};
    
    % 代数の要素
    \node at (6,0) {$\mathcal{F}_{\text{even-odd}}$ の要素：};
    \node[anchor=west] at (6,-0.8) {1. $\emptyset$};
    \node[anchor=west] at (6,-1.4) {2. $E = \{0, 2, 4\}$};
    \node[anchor=west] at (6,-2.0) {3. $E^c = \{1, 3, 5\}$};
    \node[anchor=west] at (6,-2.6) {4. $\Omega = \{0, 1, 2, 3, 4, 5\}$};
\end{tikzpicture}
\end{center}

\subsection{代数の階層}

\begin{theorem}[代数の包含関係]
\[
\mathcal{F}_{\text{trivial}} \subseteq \mathcal{F}_{\text{even-odd}} \subseteq \mathcal{F}_{\text{power}}
\]

すなわち：
\begin{itemize}
    \item $|\mathcal{F}_{\text{trivial}}| = 2$
    \item $|\mathcal{F}_{\text{even-odd}}| = 4$
    \item $|\mathcal{F}_{\text{power}}| = 2^6 = 64$
\end{itemize}
\end{theorem}

\begin{center}
\begin{tikzpicture}[
    node distance=1.5cm,
    box/.style={rectangle, draw, thick, text width=3.5cm, align=center, rounded corners, minimum height=1cm}
]
    \node[box, fill=blue!10] (trivial) {自明な代数\\{\small $|\mathcal{F}| = 2$}};
    \node[box, fill=green!10, right=of trivial] (evenodd) {偶奇代数\\{\small $|\mathcal{F}| = 4$}};
    \node[box, fill=red!10, right=of evenodd] (power) {全体の代数\\{\small $|\mathcal{F}| = 64$}};
    
    \draw[->, very thick] (trivial) -- (evenodd) node[midway, above, font=\small] {$\subseteq$};
    \draw[->, very thick] (evenodd) -- (power) node[midway, above, font=\small] {$\subseteq$};
    
    \node[below=0.5cm of trivial] {\small 最小};
    \node[below=0.5cm of power] {\small 最大};
\end{tikzpicture}
\end{center}

\section{可測関数}\label{sec:measurable}

\subsection{可測関数の定義}

代数間の「構造を保存する写像」が可測関数である。

\begin{definition}[可測関数]
可測空間 $(\Omega, \mathcal{F})$ から $(\Omega', \mathcal{G})$ への関数 $f : \Omega \to \Omega'$ が \textbf{$(\mathcal{F}, \mathcal{G})$-可測}（measurable）であるとは：
\[
\forall B \in \mathcal{G}, \quad f^{-1}(B) \in \mathcal{F}
\]
\end{definition}

\begin{insight}
\textbf{直感的理解}：

「$\Omega'$ で観測可能な事象 $B$ の逆像 $f^{-1}(B)$ が、$\Omega$ で観測可能である」

これは、「$f$ を通じて観測しても、情報が失われない」ことを意味する。
\end{insight}

\subsection{可測関数の図式}

\begin{center}
\begin{tikzpicture}[
    node distance=2.5cm,
    space/.style={rectangle, draw, thick, minimum width=2.5cm, minimum height=2cm, align=center}
]
    \node[space, fill=blue!10] (omega) {$\Omega$\\[0.3cm]{\small 可測空間}\\{\small $(\Omega, \mathcal{F})$}};
    \node[space, fill=green!10, right=of omega] (omega2) {$\Omega'$\\[0.3cm]{\small 可測空間}\\{\small $(\Omega', \mathcal{G})$}};
    
    \draw[->, very thick] (omega) -- (omega2) node[midway, above] {$f$};
    
    % 逆像の図
    \node[below=3cm of omega] (event1) {};
    \node[below=3cm of omega2] (event2) {};
    
    \draw[thick] ($ (event1) + (-1.2,0) $) rectangle ($ (event1) + (1.2,2) $);
    \draw[thick] ($ (event2) + (-1.2,0) $) rectangle ($ (event2) + (1.2,2) $);
    
    \fill[green!30] ($ (event2) + (0,0.8) $) ellipse (0.8 and 0.5);
    \node at ($ (event2) + (0,0.8) $) {$B$};
    \node at ($ (event2) + (0,1.8) $) {$\mathcal{G}$};
    
    \fill[blue!30] ($ (event1) + (0,0.8) $) ellipse (0.8 and 0.5);
    \node at ($ (event1) + (0,0.8) $) {$f^{-1}(B)$};
    \node at ($ (event1) + (0,1.8) $) {$\mathcal{F}$};
    
    \draw[<-, very thick, dashed] ($ (event1) + (0,0.8) $) -- ($ (event2) + (0,0.8) $) 
        node[midway, below] {逆像};
    
    \node at ($ (event1) + (0,-0.5) $) {\small $B \in \mathcal{G}$};
    \node at ($ (event2) + (0,-0.5) $) {\small $\Rightarrow f^{-1}(B) \in \mathcal{F}$};
\end{tikzpicture}
\end{center}

\subsection{可測関数の例}

\begin{example}[恒等関数]
恒等関数 $\mathrm{id} : \Omega \to \Omega$ は常に $(\mathcal{F}, \mathcal{F})$-可測である。

\begin{proof}
任意の $B \in \mathcal{F}$ に対して、
\[
\mathrm{id}^{-1}(B) = B \in \mathcal{F}
\]
\end{proof}
\end{example}

\begin{example}[定数関数]
定数関数 $f : \Omega \to \Omega'$, $f(x) = c$ は常に可測である。

\begin{proof}
任意の $B \in \mathcal{G}$ に対して、
\[
f^{-1}(B) = \begin{cases}
\Omega & \text{if } c \in B \\
\emptyset & \text{if } c \notin B
\end{cases}
\]

どちらも $\mathcal{F}$ に属する（$\Omega, \emptyset \in \mathcal{F}$）。
\end{proof}
\end{example}

\begin{example}[サイコロの偶奇判定]
関数 $f : \{0,1,2,3,4,5\} \to \{0,1\}$ を
\[
f(n) = \begin{cases}
0 & \text{if } n \text{ は偶数} \\
1 & \text{if } n \text{ は奇数}
\end{cases}
\]
と定義する。

このとき、$f$ は $(\mathcal{F}_{\text{even-odd}}, \mathcal{P}(\{0,1\}))$-可測である。
\end{example}

\subsection{可測関数の合成}

\begin{theorem}[可測関数の合成]
$f : \Omega \to \Omega'$ が $(\mathcal{F}, \mathcal{G})$-可測、$g : \Omega' \to \Omega''$ が $(\mathcal{G}, \mathcal{H})$-可測ならば、合成 $g \circ f : \Omega \to \Omega''$ は $(\mathcal{F}, \mathcal{H})$-可測である。
\end{theorem}

\begin{proof}
任意の $C \in \mathcal{H}$ に対して、
\begin{align*}
(g \circ f)^{-1}(C) &= f^{-1}(g^{-1}(C)) \\
&\in \mathcal{F} \quad (\because g^{-1}(C) \in \mathcal{G} \text{ かつ } f \text{ は可測})
\end{align*}
\end{proof}

\begin{center}
\begin{tikzpicture}[
    node distance=2cm,
    box/.style={rectangle, draw, thick, minimum width=1.8cm, minimum height=1.2cm, align=center}
]
    \node[box] (F) {$\mathcal{F}$};
    \node[box, right=of F] (G) {$\mathcal{G}$};
    \node[box, right=of G] (H) {$\mathcal{H}$};
    
    \draw[->, thick] (F) -- (G) node[midway, above] {$f$};
    \draw[->, thick] (G) -- (H) node[midway, above] {$g$};
    \draw[->, thick, bend right=30] (F) to node[below] {$g \circ f$} (H);
    
    \node[below=1cm of F] {\small 可測};
    \node[below=1cm of G] {\small 可測};
    \node[below=1cm of H] {\small $\Rightarrow$ 可測};
\end{tikzpicture}
\end{center}

\section{構造塔との接続}\label{sec:connection}

\subsection{フィルトレーションの概念}

\begin{definition}[フィルトレーション]
\textbf{フィルトレーション}（filtration）とは、単調増加する代数の系列：
\[
\mathcal{F}_0 \subseteq \mathcal{F}_1 \subseteq \mathcal{F}_2 \subseteq \cdots \subseteq \mathcal{F}_n
\]

各 $\mathcal{F}_t$ は時刻 $t$ で「可観測な事象の族」を表す。
\end{definition}

\begin{example}[コイン投げのフィルトレーション]
3回コイン投げ $\Omega = \{H, T\}^3$ に対して：

\begin{itemize}
    \item $\mathcal{F}_0 = \{\emptyset, \Omega\}$ (何も観測していない)
    \item $\mathcal{F}_1 = \sigma(\text{1回目の結果})$ (1回目まで観測)
    \item $\mathcal{F}_2 = \sigma(\text{1,2回目の結果})$ (2回目まで観測)
    \item $\mathcal{F}_3 = \mathcal{P}(\Omega)$ (すべて観測)
\end{itemize}

明らかに $\mathcal{F}_0 \subseteq \mathcal{F}_1 \subseteq \mathcal{F}_2 \subseteq \mathcal{F}_3$。
\end{example}

\subsection{構造塔としての定式化}

\begin{connection}
将来の \texttt{DecidableFiltration} は \texttt{StructureTowerWithMin} のインスタンスとなる：

\begin{center}
\begin{tikzpicture}[
    node distance=1.3cm,
    box/.style={rectangle, draw, thick, text width=5cm, align=center, rounded corners, minimum height=1cm}
]
    \node[box, fill=blue!10] (carrier) {carrier = FiniteAlgebra $\Omega$\\{\small (代数全体)}};
    \node[box, fill=green!10, below=of carrier] (index) {Index = $\mathbb{N}$\\{\small (時刻)}};
    \node[box, fill=yellow!10, below=of index] (layer) {layer($t$) = observableAt($t$)\\{\small (時刻 $t$ の可観測代数)}};
    \node[box, fill=red!10, below=of layer] (minlayer) {minLayer($\mathcal{F}$)\\{\small (代数 $\mathcal{F}$ が初めて現れる時刻)}};
    
    \draw[->, thick] (carrier) -- (index);
    \draw[->, thick] (index) -- (layer);
    \draw[->, thick] (layer) -- (minlayer);
    
    \node[right=0.3cm of carrier, text width=3.5cm, font=\small] {構造塔の基礎集合};
    \node[right=0.3cm of index, text width=3.5cm, font=\small] {構造塔の添字集合};
    \node[right=0.3cm of layer, text width=3.5cm, font=\small] {各層 = 各時刻の代数};
    \node[right=0.3cm of minlayer, text width=3.5cm, font=\small] {\color{red}構造塔の鍵！};
\end{tikzpicture}
\end{center}
\end{connection}

\subsection{構造塔の対応表}

\begin{center}
\begin{tabular}{|l|l|}
\hline
\textbf{構造塔の概念} & \textbf{フィルトレーションの概念} \\
\hline
carrier & 代数の全体 \\
Index & 時刻 $\mathbb{N}$ \\
layer($t$) & 時刻 $t$ の可観測代数 $\mathcal{F}_t$ \\
monotone\_level & 情報の単調増加 $\mathcal{F}_s \subseteq \mathcal{F}_t$ \\
minLayer($\mathcal{F}$) & $\mathcal{F}$ が初めて観測可能になる時刻 \\
\hline
\end{tabular}
\end{center}

\subsection{停止時間との対応}

\begin{definition}[停止時間（予告）]
ランダム変数 $\tau : \Omega \to \mathbb{N}$ が\textbf{停止時間}（stopping time）であるとは：
\[
\forall t \in \mathbb{N}, \quad \{\tau \leq t\} \in \mathcal{F}_t
\]
\end{definition}

\begin{connection}
\textbf{構造塔の視点}：

停止時間の条件「$\{\tau \leq t\} \in \mathcal{F}_t$」は、
\[
\mathrm{minLayer}(\{\tau \leq t\}) \leq t
\]
と書き直せる。

すなわち、事象 $\{\tau \leq t\}$ が「時刻 $t$ までに初めて観測可能になる」。
\end{connection}

\begin{example}[初めて表が出る時刻]
コイン投げで「初めて表が出る時刻」$\tau$ は停止時間である。

なぜなら、時刻 $t$ までに「$\tau \leq t$ か？」（= $t$ 回目までに少なくとも1回表が出たか？）は、時刻 $t$ の情報 $\mathcal{F}_t$ だけで判定できるから。
\end{example}

\section{まとめと展望}\label{sec:conclusion}

\subsection{本ファイルの成果}

\texttt{DecidableAlgebra.lean} では、以下を達成した：

\begin{enumerate}
    \item \textbf{有限代数の定義}
    \begin{itemize}
        \item Boolean 演算で閉じた事象族
        \item $\sigma$-代数の有限版
    \end{itemize}
    
    \item \textbf{基本的な性質}
    \begin{itemize}
        \item 積・差で閉じている
        \item 有限和・有限積で閉じている
    \end{itemize}
    
    \item \textbf{具体例}
    \begin{itemize}
        \item 自明な代数（最小）
        \item 全体の代数（最大）
        \item 偶奇代数（サイコロの例）
    \end{itemize}
    
    \item \textbf{可測関数}
    \begin{itemize}
        \item 代数間の構造を保存する写像
        \item 恒等関数と合成
    \end{itemize}
    
    \item \textbf{構造塔との接続}
    \begin{itemize}
        \item フィルトレーションの準備
        \item minLayer による特徴づけの予告
    \end{itemize}
\end{enumerate}

\subsection{階層構造の全体像}

\begin{center}
\begin{tikzpicture}[
    node distance=1.5cm,
    level/.style={rectangle, draw, thick, minimum width=6.5cm, minimum height=1cm, align=center},
    arrow/.style={->, thick}
]
    \node[level, fill=green!10] (level1) {DecidableEvents.lean\\{\small 事象の decidability}};
    \node[level, fill=blue!10, below=of level1] (level2) {DecidableAlgebra.lean\\{\small Boolean 代数・可測性}};
    \node[level, fill=yellow!10, below=of level2] (level3) {DecidableFiltration.lean\\{\small StructureTowerWithMin}};
    \node[level, fill=orange!10, below=of level3] (level4) {ComputableStoppingTime.lean\\{\small 停止時間 = minLayer}};
    \node[level, fill=red!10, below=of level4] (level5) {AlgorithmicMartingale.lean\\{\small Optional Stopping Theorem}};
    
    \draw[arrow] (level1) -- (level2) node[midway, right, font=\small] {構造化};
    \draw[arrow] (level2) -- (level3) node[midway, right, font=\small] {時間構造};
    \draw[arrow] (level3) -- (level4) node[midway, right, font=\small] {決定理論};
    \draw[arrow] (level4) -- (level5) node[midway, right, font=\small] {マルチンゲール};
    
    \node[left=0.3cm of level2, text width=2cm, font=\small] {\color{blue}今回完了};
\end{tikzpicture}
\end{center}

\subsection{今後の課題}

\begin{enumerate}
    \item \textbf{DecidableFiltration.lean}（次のファイル）
    \begin{itemize}
        \item StructureTowerWithMin のインスタンス化
        \item 単調増加する代数の系列
        \item minLayer = 初めて可観測になる時刻
    \end{itemize}
    
    \item \textbf{ComputableStoppingTime.lean}
    \begin{itemize}
        \item 停止時間の定義
        \item minLayer による特徴づけ
        \item 具体的な停止時間の計算
    \end{itemize}
    
    \item \textbf{AlgorithmicMartingale.lean}
    \begin{itemize}
        \item マルチンゲールの定義
        \item Optional Stopping Theorem
        \item 有限計算による証明
    \end{itemize}
\end{enumerate}

\subsection{教育的価値}

\begin{keypoint}
本ファイルは、以下の重要な洞察を提供する：

\begin{enumerate}
    \item \textbf{代数的構造の重要性}：事象の族も代数的構造を持つ
    \item \textbf{情報の段階的増加}：フィルトレーションの直感
    \item \textbf{構造塔との自然な対応}：layer = 代数、minLayer = 初出時刻
    \item \textbf{可測性の本質}：「観測可能性」の形式化
\end{enumerate}
\end{keypoint}

\subsection{Lean 4 形式化の意義}

\begin{insight}
型システムによる保証：

\begin{itemize}
    \item \textbf{代数の公理}：補集合・和で閉じていることを型で保証
    \item \textbf{可測性}：逆像が代数に属することを証明で保証
    \item \textbf{合成の可測性}：可測関数の合成が自動的に可測
\end{itemize}

これにより、「可測だと思ったが実は違った」という状況を防ぐ。
\end{insight}

\section*{参考文献}

\begin{enumerate}
    \item Billingsley, Patrick. \textit{Probability and Measure}. 3rd ed. Wiley, 1995.
    \item Williams, David. \textit{Probability with Martingales}. Cambridge University Press, 1991.
    \item Kallenberg, Olav. \textit{Foundations of Modern Probability}. 2nd ed. Springer, 2002.
    \item The mathlib Community. \textit{The Lean Mathematical Library}. \url{https://github.com/leanprover-community/mathlib4}
    \item 本プロジェクト：\texttt{DecidableEvents.lean}, \texttt{Bourbaki\_Lean\_Guide.lean}, \texttt{CAT2\_complete.lean}
\end{enumerate}

\section*{付録 A：Lean コードへのリンク}

\subsection*{主要定義}

\begin{itemize}
    \item \textbf{FiniteAlgebra}：170--179 行目
    \item \textbf{has\_univ}（全事象の包含）：196--199 行目
    \item \textbf{closed\_intersection}（積で閉じている）：206--217 行目
    \item \textbf{closed\_finite\_union}（有限和で閉じている）：261--273 行目
\end{itemize}

\subsection*{具体例}

\begin{itemize}
    \item \textbf{trivial}（自明な代数）：330--348 行目
    \item \textbf{powerSet}（全体の代数）：358--362 行目
    \item \textbf{evenOddAlgebra}（偶奇代数）：430--464 行目
\end{itemize}

\subsection*{可測関数}

\begin{itemize}
    \item \textbf{MeasurableFunction}：518--524 行目
    \item \textbf{id}（恒等可測関数）：534--539 行目
    \item \textbf{comp}（可測関数の合成）：542--548 行目
\end{itemize}

\subsection*{構造塔との接続}

\begin{itemize}
    \item \textbf{フィルトレーションの説明}：562--622 行目
\end{itemize}

\end{document}
